\chapter{Low-Level Control and Sensor Fusion}
\label{ch:commande_bas_niveau}

\section{Motor Control (\texttt{motor\_bridge\_node})}
\label{sec:cbn_motor_bridge}

This node (package \texttt{esibot\_hardware}) is the core of the robot's real-time
control. On the Raspberry\,Pi\,4, it plays the role originally assigned to the
ESP32/\texttt{esibot\_driver} pair in the specification: receiving \topic{/cmd\_vel},
wheel velocity control, encoder reading, and publishing raw odometry.

\subsection{Control Loop}
\label{sec:cbn_loop}

The node runs a periodic loop at \texttt{control\_hz} (10\,Hz by default) which, on each
cycle:
\begin{enumerate}
	\item reads the pulse counts accumulated since the last cycle for each wheel;
	\item computes the measured rotation speed in \unit{\rpm} (equation~\ref{eq:rpm_measure});
	\item applies a \gls{pid} controller (one per wheel) to bring the measured speed closer to
	      the setpoint derived from the differential kinematics (equation~\ref{eq:diff_drive});
	\item converts the controller output into a PWM duty cycle via \texttt{pwm\_mapper.py};
	\item applies an optional heading correction based on the \gls{imu} (yaw control);
	\item publishes the integrated odometry on \topic{/odom} and the counter status on the
	      diagnostic topics (\topic{/encoder\_ticks\allowbreak/left}, \topic{/encoder\_ticks\allowbreak/right}).
\end{enumerate}

\subsection{Per-Wheel \gls{pid} Controller}
\label{sec:cbn_pid}

Each wheel has its own independent \gls{pid} controller, with default gains $K_p$,
$K_i$, $K_d$ configurable via the \ros{} parameter server (and adjustable during
operation without recompilation). The controller output is saturated, and an
anti-\emph{windup} mechanism limits accumulation of the integral term when the command is
saturated, avoiding overshoot during abrupt setpoint changes.

\begin{notebox}[title={Safety Loop}]
	The node implements a \emph{watchdog} timer: if no \topic{/cmd\_vel} message is
	received for a duration \texttt{cmd\_vel\_timeout} (typically a few hundred
	milliseconds), both motors are automatically cut off. This mechanism protects the
	robot in case of network communication loss with the teleoperation node or Nav2.
\end{notebox}

\section{Wheel Speed to PWM Conversion}
\label{sec:cbn_pwm_mapper}

The \texttt{pwm\_mapper.py} module converts a desired linear wheel speed (in
\unit{\metre\per\second}) into a normalized PWM duty cycle. In \texttt{working\_build},
this conversion is no longer a simple theoretical calibration: it loads the
\texttt{\textasciitilde/.esibot/pwm\_rpm\_map.json} file, generated on the robot, and
applies a separate affine regression for each wheel and direction of travel. The
conversion combines:
\begin{itemize}
	\item a mechanical dead zone (\texttt{min\_duty\_cycle}, 15\,\% by default and 22\,\%
	      in some launches without calibration);
	\item \texttt{a}/\texttt{b} coefficients calibrated for \texttt{left/right} and
	      \texttt{forward/reverse}, to compensate for efficiency differences between the
	      two motors;
	\item a cap at \texttt{max\_pwm}/\texttt{max\_rpm} derived from calibration, which
	      avoids commanding into a saturation or mechanical-noise region.
\end{itemize}

\section{Calibration Tool (\texttt{motor\_calibrator.py})}
\label{sec:cbn_calibrator}

A standalone calibration script (outside the \ros{} graph, run from the command line)
empirically characterizes the relationship between PWM duty cycle and actual rotation
speed (in \unit{\rpm}) for each motor. It sweeps duty cycles from 15\,\% to 100\,\% in
5\,\% steps, measures two one-second samples per point, keeps the median, then stops the
sweep as soon as an RPM drop indicates saturation or an unstable reading. The result is
saved to \texttt{\textasciitilde/.esibot/pwm\_rpm\_map.json} along with the wheel radius
used (0.033\,m), the PWM frequency (500\,Hz) and the forward/reverse coefficients for
each wheel. This tool is essential for limiting motor mechanical asymmetry, but it alone
is not sufficient to guarantee correct odometry if the actual wheel radius, wheelbase and
\gls{slam} parameters are not recalibrated together (Chapter~\ref{ch:limitations}).

\section{Inertial Measurement Unit and Heading Stabilization}
\label{sec:cbn_imu}

The \texttt{mpu6050\_node} node (package \texttt{esibot\_sensors}) starts with a
stationary gyroscope bias calibration over 100 samples, then continuously publishes
(50\,Hz):
\begin{itemize}
	\item \topic{/imu/data\_raw} (\texttt{sensor\_msgs/Imu}) --- raw linear accelerations
	      $(a_x, a_y, a_z)$ and angular velocities $(\omega_x, \omega_y, \omega_z)$;
	\item \topic{/imu/temperature} (\texttt{sensor\_msgs/Temperature}) --- internal sensor
	      temperature, useful for thermal diagnostics.
\end{itemize}

\texttt{motor\_bridge\_node} subscribes to \topic{/imu/data\_raw} to apply closed-loop
yaw correction when the robot is moving forward and no deliberate rotation is commanded.
The controller compensates for residual yaw rate measured by the \gls{imu}, with a dead
band of 0.02\,rad/s and a maximum correction of 40\,RPM per wheel.

\section{Data Fusion (EKF Filter)}
\label{sec:cbn_ekf}

Raw odometry (\topic{/odom}, from encoders alone) and inertial data
(\topic{/imu/data\_raw}) are fused by an \gls{ekf} provided by the standard
\texttt{robot\_localization} package (\texttt{ekf\_filter\_node} node), configured in
two-dimensional mode (\texttt{two\_d\_mode: true}). The \texttt{working\_build}
\texttt{ekf\_odom\_imu.yaml} file is deliberately conservative: it fuses the odometry
velocities $(v_x,v_y,\dot\theta)$ and only the yaw rate from the \gls{imu}. It injects
neither absolute orientation nor linear acceleration, since the MPU6050 does not publish
an estimated orientation and its raw acceleration contains gravity and residual bias.
The \topic{/odometry/filtered} output thus provides a more stable local odometry for
\gls{slam}/Nav2, while \texttt{slam\_toolbox} or \gls{amcl} retain responsibility for the
global \topic{map}~$\rightarrow$~\topic{odom} correction.
